% !TEX root = paper.tex
\begin{abstract}
Diffusion Transformers (DiT)
have emerged as a leading paradigm
for high-fidelity image and video generation,
demonstrating remarkable generative capabilities.
However,
their practical deployment
remains constrained by the inherent computational burden
of iterative denoising,
substantially limiting their applicability
in latency-sensitive and real-time scenarios.
% While prior work has explored cache-based mechanisms
% to amortize computational costs
% through selective feature reuse across timesteps,
% existing methods
% typically rely on static reuse schedules and heuristic thresholds.
Because inter-timestep coherence
varies substantially over the denoising trajectory,
fixed policies are brittle:
they either skip essential computations when activations diverge,
incurring noticeable fidelity degradation,
or remain overly conservative,
leaving much of the redundancy unexploited.
This creates a persistent tension
between computational efficiency and generation quality.
To address this challenge
of inference latency and generation fidelity,
we propose $D^3$-Cache,
a training-free diffusion model acceleration method
that leverages low-rank linear dynamic mode decomposition
to forecast transformer output tensors,
thereby enabling omission
of redundant forward computations.
This approach allows us to extrapolate transformer outputs
with minimal quality degradation.
Extensive empirical validation
across diverse image and video generation tasks
demonstrates the efficacy of our method.
$D^3$-Cache achieves $2.49-4.32\times$ inference acceleration
on state-of-the-art models
while preserving generation fidelity and quality.
$D^3$-Cache is open-source and available
at \url{https://anonymous.4open.science/r/dmd-E5A2/}.
\end{abstract}
